% SGA 3, Expose VI_B, Loop 1 English reconstruction.
% Scope: Section 11 heading, paragraph 11.0, Notation,
% Lemma 11.1, Corollary 11.2, Definition 11.3, and
% Proposition 11.3.1, including editor's notes 104--107.
% Authority: Polo--Gille Exp6B-13oct24.pdf, component-local
% pages 76--78 (combined-reader pages 417--419; re-edition
% headers 406--408; original-margin printed pages 396--398).
% Paragraph 11.4 begins on local page 78 and is excluded.
% OCR and Reinhold English Markdown are comparison witnesses only.
% Loop 2 replaces both source-bound commutative-square crops with
% faithful native LaTeX reconstructions; both are closed.

\section*{11. 仿射群概形}
\phantomsection\label{sec:sga3-VIB-11}%
\addcontentsline{toc}{section}{11. 仿射群概形}

\medskip
\phantomsection\label{par:sga3-VIB-11.0}%
\noindent\textbf{11.0. 回顾.---}%
\footnote{\label{note:sga3-VIB-104}编者注: 此回顾段落为新增内容.}
设 \(q:X\to S\) 为拟紧且拟分离的概形态射
\textup{(参见 EGA~IV\(_1\), 1.1 和 1.2)},
\(\mathcal F\) 为拟凝聚 \(\mathcal O_X\)-模.
回顾, \(q_*(\mathcal F)\) 为拟凝聚 \(\mathcal O_S\)-模
\textup{(EGA~I, 9.2.1)}.
此外, 由 EGA~III, 1.4.15 \textup{(由 EGA~IV\(_1\), 1.7.21 补全)},
有下述\hyperref[concl:sga3-VIB-11.0-c]{\textup{(c)}},
而\emph{同上}的证明也给出\hyperref[concl:sga3-VIB-11.0-a]{\textup{(a)}}
和\hyperref[concl:sga3-VIB-11.0-b]{\textup{(b)}}:
\begin{enumerate}
\item[\textup{(a)}]\phantomsection\label{concl:sga3-VIB-11.0-a}
  若 \(\mathcal F\) 为拟凝聚子模 \(\mathcal F_\alpha\) 的滤过归纳极限, 则
  \[
    q_*(\mathcal F)
      =\varinjlim_\alpha q_*(\mathcal F_\alpha).
  \]

\item[\textup{\SGAInlineRefTarget{par:sga3-VIB-11.0:item:b}(b)}]\phantomsection\label{concl:sga3-VIB-11.0-b}
  若 \(\mathcal E\) 为平坦 \(\mathcal O_S\)-模, 则典范态射
  \phantomsection\label{form:sga3-VIB-11.0-projection}%
  \[
    \mathcal E\otimes_{\mathcal O_S}q_*(\mathcal O_X)
      \longrightarrow q_*q^*(\mathcal E)
  \]
  为同构.

\item[\textup{(c)}]\phantomsection\label{concl:sga3-VIB-11.0-c}
  设 \(p:S'\to S\) 为平坦态射, \(q':X'\to S'\) 由 \(q\) 经基变换得到,
\(\mathcal F'\) 为 \(\mathcal F\) 在 \(X'\) 上的逆像. 则典范态射
  \phantomsection\label{form:sga3-VIB-11.0-flat-base-change}%
  \[
    p^*q_*(\mathcal F)\longrightarrow q'_*(\mathcal F')
  \]
  为同构.
\end{enumerate}

事实上, 设 \(U=\operatorname{Spec}(A)\) 为 \(S\) 的任意仿射开子集.
由假设, \(q^{-1}(U)\) 为如下仿射开子集之并:
\[
  V_i=\operatorname{Spec}(B_i),\qquad i=1,\ldots,n,
\]
且每个交 \(V_i\cap V_j\) 为有限多个仿射开子集
\[
  W_{ijk}=\operatorname{Spec}(C_{ijk}).
\]
之并. 于是 \(\Gamma(U,q_*(\mathcal F))
  =\Gamma(q^{-1}(U),\mathcal F)\) 为下述态射的核:
\phantomsection\label{form:sga3-VIB-11.0-Cech-F}%
\[
  \bigoplus_{i=1}^{n}\Gamma(V_i,\mathcal F)
    \longrightarrow
  \bigoplus_{i,j,k}\Gamma(W_{ijk},\mathcal F).
\]
由此得到断言~\hyperref[concl:sga3-VIB-11.0-a]{\textup{(a)}},
因为上面的每一项都与滤过归纳极限可交换
\textup{(\(V_i\) 和 \(W_{ijk}\) 仿射, 从而拟紧)}.

下面证明\hyperref[concl:sga3-VIB-11.0-b]{\textup{(b)}}.
模 \(E=\Gamma(U,\mathcal E)\) 在 \(A\) 上平坦,
且 \(\Gamma(U,q_*q^*(\mathcal E))\) 是下述态射的核 \(K(E)\):
\phantomsection\label{form:sga3-VIB-11.0-Cech-E}%
\[
  \bigoplus_{i=1}^{n}B_i\otimes_AE
    \longrightarrow
  \bigoplus_{i,j,k}C_{ijk}\otimes_AE.
\]
由于 \(E\) 在 \(A\) 上平坦, 此核可等同于
\[
  K(A)\otimes_AE
    =\mathcal O_X(q^{-1}(U))\otimes_AE.
\]
最后, 若 \(U'\) 为 \(S'\) 中位于 \(U\) 之上的任意仿射开子集,
则 \(A'=\mathcal O(U')\) 为平坦 \(A\)-代数, 相同的论证给出%
\begingroup
\renewcommand{\thefootnote}{\fnsymbol{footnote}}%
\footnote[1]{\label{note:sga3-VIB-11.0-U-prime}
英译者注: 法文重版将下式右端印作 \(q'^{-1}(U)\).
由于 \(q':X'\to S'\), 而刚引入的仿射开子集为 \(U'\subset S'\),
故逆像为 \(q'^{-1}(U')\).}%
\endgroup
\phantomsection\label{form:sga3-VIB-11.0-local-base-change}%
\[
  \mathcal F(q^{-1}(U))\otimes_AA'
    \xrightarrow{\ \sim\ }
  \mathcal F'(q'^{-1}(U')).
\]

\medskip
\phantomsection\label{not:sga3-VIB-11-A}%
\noindent\textbf{记号.---}
设 \(S\) 为概形, \(X\) 为 \(S\)-概形,
\(f:X\to S\) 为结构态射. 令
\phantomsection\label{form:sga3-VIB-11-A-X}%
\[
  \mathcal A(X)=f_*(\mathcal O_X).
\]

\medskip
\phantomsection\label{lem:sga3-VIB-11.1}%
\noindent\textbf{引理 11.1.---}
\emph{设 \(X\) 和 \(Y\) 为两个拟紧且拟分离的 \(S\)-概形,
\(f:X\to S\) 和 \(g:Y\to S\) 为其结构态射. 则典范同态}
\phantomsection\label{form:sga3-VIB-11.1-phi}%
\[
  \varphi:
  \mathcal A(X)\otimes_{\mathcal O_S}\mathcal A(Y)
    \longrightarrow
  \mathcal A(X\times_SY)
\]
\emph{在下列各情形中为同构:}%
\footnote{\label{note:sga3-VIB-105}编者注: 若 \(k\) 为域, \(X\) 为无限多个
\(S=\operatorname{Spec}k\) 的副本的不交并
\textup{(因而 \(X\) 不拟紧)}, 则 \(\mathcal A(X)=k^X\), 且典范态射
\[
  k^X\otimes_k k^X\longrightarrow k^{X\times X}
\]
不为满射.}
\begin{enumerate}
\item[\textup{(a)}]\phantomsection\label{concl:sga3-VIB-11.1-a}
  \emph{\(f\) 和 \(g\) 仿射;}

\item[\textup{(b)}]\phantomsection\label{concl:sga3-VIB-11.1-b}
  \emph{\(g\) \textup{(或 \(f\))} 平坦且仿射;}

\item[\textup{(c)}]\phantomsection\label{concl:sga3-VIB-11.1-c}
  \emph{\(g\) 平坦, 且 \(f_*(\mathcal O_X)\) 为平坦 \(\mathcal O_S\)-模.}
\end{enumerate}

在情形~\hyperref[concl:sga3-VIB-11.1-b]{\textup{(b)}} 中,
假设平坦且仿射的是 \(g\). 令
\[
  S'=\operatorname{Spec}\mathcal A(X),\qquad
  Y'=Y\times_SS',\qquad
  g'=g\times_SS',
\]
并用 \(v:S'\to S\) 表示结构态射:
\begin{center}
  \phantomsection\label{diag:sga3-VIB-11.1-base-change}%
  \begin{tikzcd}[row sep=large,column sep=large]
    Y' \arrow[r] \arrow[d,"g'"'] &
    Y \arrow[d,"g"] \\
    \operatorname{Spec}\mathcal A(X)=S'
      \arrow[r,"v"'] &
    S
  \end{tikzcd}
\end{center}
在情形~\hyperref[concl:sga3-VIB-11.1-a]{\textup{(a)}}
和\hyperref[concl:sga3-VIB-11.1-b]{\textup{(b)}} 中, 态射 \(g\) 仿射,
故由 EGA~II, 1.5.2 得
\phantomsection\label{form:sga3-VIB-11.1-eq1}%
\[
  \tag{1}
  g'_*(\mathcal O_{Y'})
    =v^*g_*(\mathcal O_Y)
    =\mathcal A(Y)\otimes_{\mathcal O_S}\mathcal O_{S'}.
\]
在情形~\hyperref[concl:sga3-VIB-11.1-c]{\textup{(c)}} 中,
由\hyperref[concl:sga3-VIB-11.0-c]{11.0\textup{(c)}} 也有同一等式,
因为 \(S'\) 在 \(S\) 上平坦, 且 \(g\) 拟紧且拟分离.

另一方面, 由 EGA~II, 1.2.7, 态射 \(f:X\to S\)
通过态射 \(p:X\to S'\) 经由 \(v\) 分解, 且
\[
  p_*(\mathcal O_X)=\mathcal O_{S'},
  \qquad
  X\times_SY=X\times_{S'}Y'.
\]
由于 \(f\) 拟分离, \(p\) 也拟分离
\textup{(EGA~IV\(_1\), 1.2.2)};
又由于 \(f\) 拟紧且 \(v\) 拟分离, \(p\) 也拟紧
\textup{(EGA~IV\(_1\), 1.2.4)}. 考虑笛卡儿方图
\begin{center}
  \phantomsection\label{diag:sga3-VIB-11.1-cartesian}%
  \begin{tikzcd}[row sep=large,column sep=large]
    X\times_SY \arrow[r,"p'"] \arrow[d] &
    Y' \arrow[d,"g'"] \\
    X \arrow[r,"p"'] &
    S'
  \end{tikzcd}
\end{center}
在情形~\hyperref[concl:sga3-VIB-11.1-b]{\textup{(b)}}
和\hyperref[concl:sga3-VIB-11.1-c]{\textup{(c)}} 中,
概形 \(Y\) 在 \(S\) 上平坦, 因而 \(Y'\) 在 \(S'\) 上平坦.
再次应用\hyperref[concl:sga3-VIB-11.0-c]{11.0\textup{(c)}}, 得到
\phantomsection\label{form:sga3-VIB-11.1-eq2}%
\[
  \tag{2}
  p'_*(\mathcal O_{X\times_SY})
    =g'^*p_*(\mathcal O_X)
    =g'^*(\mathcal O_{S'})
    =\mathcal O_{Y'}.
\]
在情形~\hyperref[concl:sga3-VIB-11.1-a]{\textup{(a)}} 中也有同一等式,
因为此时 \(p\) 和 \(p'\) 为同构.

最后, 由于 \(v\) 仿射, EGA~II, 1.4.7 给出
\[
  v_*(\mathcal F\otimes_{\mathcal O_S}\mathcal O_{S'})
    =\mathcal F\otimes_{\mathcal O_S}\mathcal A(X)
\]
其中 \(\mathcal F\) 为任意拟凝聚 \(\mathcal O_S\)-模.
结合\hyperref[form:sga3-VIB-11.1-eq2]{\textup{(2)}}
和\hyperref[form:sga3-VIB-11.1-eq1]{\textup{(1)}}, 得到
\[
\begin{aligned}
  \mathcal A(X\times_SY)
    &=v_*g'_*p'_*(\mathcal O_{X\times_SY})\\
    &\overset{\textup{(2)}}{=}v_*g'_*(\mathcal O_{Y'})
     \overset{\textup{(1)}}{=}
       v_*(\mathcal A(Y)\otimes_{\mathcal O_S}\mathcal O_{S'})\\
    &=\mathcal A(Y)\otimes_{\mathcal O_S}\mathcal A(X).
\end{aligned}
\]

\medskip
\phantomsection\label{cor:sga3-VIB-11.2}%
\noindent\textbf{推论 11.2.---}
\emph{函子}
\[
  X\longmapsto\operatorname{Spec}\mathcal A(X),
\]
\emph{从 \((\mathbf{Sch}/S)\) 的由那些在 \(S\) 上平坦、拟紧且拟分离,
并使 \(\mathcal A(X)\) 为平坦 \(\mathcal O_S\)-模的
\(S\)-概形 \(X\) 组成的满子范畴, 到平坦仿射
\(S\)-概形范畴, 与有限积可交换, 因而将 \(S\)-群映为 \(S\)-群.}

\medskip
\phantomsection\label{def:sga3-VIB-11.3}%
\noindent\textbf{定义 11.3.---}
\emph{给定在 \(S\) 上拟紧且拟分离的平坦 \(S\)-群 \(G\),
使 \(\mathcal A(G)\) 在 \(\mathcal O_S\) 上平坦,}%
\footnote{\label{note:sga3-VIB-106}编者注: 若 \(S\) 为维数不超过 \(2\) 的正则局部 Noether 概形,
\(X\) 为平坦、拟紧且拟分离的 \(S\)-概形,
则 \(\mathcal A(X)\) 为平坦 \(\mathcal O_S\)-模;
参见 \textup{[Ray70a]}, VII, \SGARefLink{sga3:VIIB:3.2:prop}{3.2}.}
\emph{将 \(S\)-群}
\[
  G_{\mathrm{af}}=\operatorname{Spec}\mathcal A(G)
\]
\emph{记为 \(G_{\mathrm{af}}\), 称为 \(G\) 的\emph{仿射包络 (affine envelope)}.}

\medskip
\phantomsection\label{prop:sga3-VIB-11.3.1}%
\noindent\textbf{命题 11.3.1.---}%
\footnote{\label{note:sga3-VIB-107}编者注: 此命题为新增内容;
关于态射 \(G\to G_{\mathrm{af}}\) 及其核的研究,
亦见下文增补的\hyperref[sec:sga3-VIB-12]{第~12~节}.}
\emph{典范态射}
\[
  \tau_G:G\longrightarrow G_{\mathrm{af}}
\]
\emph{为 \(S\)-群的态射. 此外, 它有下述普遍性质:}
\begin{enumerate}
\item[\textup{(i)}]\phantomsection\label{concl:sga3-VIB-11.3.1-i}
  \emph{对每个 \(S\)-概形态射 \(\varphi:G\to H\),
其中 \(H\) 在 \(S\) 上仿射, 存在唯一的 \(S\)-概形态射
\(\varphi':G_{\mathrm{af}}\to H\), 使得}
  \[
    \varphi=\varphi'\circ\tau_G.
  \]

\item[\textup{(ii)}]\phantomsection\label{concl:sga3-VIB-11.3.1-ii}
  \emph{此外, 若 \(H\) 为 \(S\)-群, 且 \(\varphi\) 为 \(S\)-群的态射,
则 \(\varphi'\) 也为群态射.}
\end{enumerate}
